\section{Background}

% ============================================================
% 2.1  Accountability in Multi-Agent AI Systems
% ============================================================
Accountability is a foundational requirement for any AI system that
operates autonomously in high-stakes environments.  In multi-agent
architectures, accountability is not merely a property of individual
agents but an emergent property of the interaction protocols,
authority hierarchies, and logging mechanisms that bind agents into a
coherent system.  When an autonomous agent takes an action that
produces a harmful outcome, it must be possible to trace that outcome
back through the chain of decisions, communications, and dependencies
that gave rise to it \cite{dijkstra1982}.  This traceability is the
core of accountability: it closes the gap between action and consequence
in systems where human oversight is delayed, distributed, or absent.

The challenge of accountability in multi-agent AI is compounded by the
semantic opacity that emerges when multiple agents share goals but
pursue them through opaque inference chains.  Unlike in traditional
software systems where a call stack or transaction log provides a
faithful record of execution, a multi-agent system may produce
decisions that cannot be explained by inspecting a single agent's
internal state.  This has led a growing body of work to emphasize the
necessity of structured audit trails that are co-designed with the
agent architecture itself, rather than retrofitted as an after-thought
\cite{bansal2019}.  Effective accountability mechanisms must record not
only the final decision of an agent but the inputs it received, the
alternatives it considered, and the reasoning it applied—regardless of
whether that reasoning is symbolically expressible.  In short,
accountability is a design discipline, not a compliance afterthought.

% ============================================================
% 2.2  Human-in-the-Loop Design
% ============================================================
Human-in-the-loop (HITL) design embeds human authority at one or more
decision boundaries within an autonomous system, ensuring that a human
retains meaningful oversight of actions that carry material risk.  The
central insight of HITL is that automation does not eliminate human
responsibility; it redefines it—shifting the human from an active
operator to a supervisory referee whose intervention is both possible
and consequential \cite{wiener1948}.  When designed correctly, HITL
preserves the throughput advantages of automation while preserving the
capacity for moral reasoning, contextual judgment, and recourse that
only humans can provide.

Two dimensions of HITL design are particularly relevant to the
Bailout Protocol.  The first is the placement of decision authority:
at what stage of a cascade should a human be empowered to interrupt?
Prior work has shown that the effectiveness of a human interrupt
diminishes rapidly once an autonomous process has entered a
high-rate execution phase, a phenomenon sometimes described as
\emph{automation bias} \cite{bansal2019}.  Humans who have delegated
routine decisions to an autonomous system tend to approve rather than
scrutinize, particularly under time pressure.  This suggests that HITL
placement must be tied not to the pace of the machine but to the
consequence class of the action being contemplated.  The second
dimension is the fidelity of the human's situational awareness at the
moment of intervention.  A human who is consulted after a cascade has
already propagated through multiple agents cannot meaningfully assess
whether any individual decision was sound, because the original
contexts have been obscured by subsequent processing
\cite{dijkstra1982}.  HITL mechanisms therefore require rich,
concise, and human-interpretable state summaries—designed before the
cascade, not improvised during it.

% ============================================================
% 2.3  Fail-Safe Design in Safety-Critical Systems
% ============================================================
The concept of a fail-safe system predates modern computing by
decades.  A fail-safe design is one in which the system degrades
gracefully and predictably when stressed beyond its operational
envelope, producing outcomes that are less harmful than those that
would result from continued unattended operation \cite{tristr81}.
Classically, fail-safe engineering distinguishes between systems that
fail gracefully (degrading to a lower functionality mode) and systems
that fail stop (halting completely upon detecting an anomalous
condition).  The choice between these modes depends on the cost of
the halt versus the cost of continued operation in a compromised state.

In the context of autonomous agents, fail-safe design takes on new
complexity.  An agent operating in a dynamic environment may encounter
conditions that were not anticipated at design time, meaning that a
static fail-safe threshold is insufficient \cite{bansal2019}.  Rather,
a robust fail-safe architecture must include mechanisms for the agent
to recognize that it has entered an unanticipated state, to reduce its
autonomy in a principled manner, and to communicate its uncertainty
to external oversight before taking irreversible actions.  Dijkstra's
observation that programs must be designed to behave correctly not only
in their intended domain but in the face of all foreseeable errors
remains the foundational principle \cite{dijkstra1982}.  The novel
challenge for agentic systems is that the space of foreseeable errors
is orders of magnitude larger than in traditional software, and the
errors that arise may be emergent properties of agent interaction
rather than failures of any individual component.

Modern safety-critical systems, particularly those deployed in
transportation and industrial automation, have converged on a
principle of layered defense-in-depth \cite{tristr81}, in which
multiple independent mechanisms are layered such that the failure of
any single layer does not result in a catastrophic loss of safety.
The Bailout Protocol draws on this tradition, treating the bailouts
not as a single defensive layer but as part of a broader architecture
in which each layer—task execution, oversight, and governance—can
independently detect anomalies and initiate protective action.

% ============================================================
% 2.4  The BX3 Framework: Three-Layer Model
% ============================================================
The BX3 Framework organizes an agentic system into three concentric
layers of increasing abstraction and decreasing autonomy: the
\textbf{Task Layer}, the \textbf{Oversight Layer}, and the
\textbf{Governance Layer}.  The Task Layer is the innermost and most
autonomous: it performs the operational work of the system—planning,
reasoning, and executing actions in the world.  It is analogous to the
runtime kernel of a traditional operating system and corresponds to
what Wiener described as the \emph{executive function} of an automated
system \cite{wiener1948}.

The Oversight Layer sits above the Task Layer and is responsible for
monitoring the quality, safety, and alignment of decisions made at the
Task Layer.  It implements the human-in-the-loop interfaces, audit
logging, and anomaly detection mechanisms described in Sections 2.2
and 2.3.  Critically, the Oversight Layer is itself an agentic system:
it reasons about the behavior of the Task Layer and can initiate
corrective actions without waiting for a human interrupt.  This
self-referential structure—agents monitoring agents—is necessary
because the latency of human supervision is incompatible with the pace
at which modern agentic systems operate \cite{bansal2019}.  The
Oversight Layer is thus the primary locus of accountability in the BX3
Framework.

The Governance Layer is the outermost and slowest-moving layer.  It
encodes the terminal values, policy constraints, and bailout triggers
that govern the behavior of the system as a whole.  Unlike the Task
Layer and Oversight Layer, which are designed to operate at runtime,
the Governance Layer is primarily a design-time artifact: it is the
set of constraints that the system cannot recover from autonomously.
When the Oversight Layer detects a condition that violates aGovernance
Layer constraint, it must escalate—not to a more capable agent, but
to a human.  This is the structural basis of the Bailout Protocol:
the point at which the system exhausts its internal recovery options
and transfers authority back to human governance \cite{tristr81}.

The three-layer model provides a principled answer to the question of
where to place accountability: it is distributed across the layers
according to the nature of the decision.  Operational decisions are
accountable to the Oversight Layer; safety and alignment decisions
are accountable to the Governance Layer.  The Bailout Protocol is
the formal mechanism by which the Governance Layer exercises this
authority—declaring that a condition has arisen that requires human
intervention, halting the autonomous cascade, and preserving the
state necessary for the human to reconstruct, assess, and resolve
the situation.  This architecture ensures that fail-safe design is
not an optional enhancement but a structural property of the system,
enforced by the separation of concerns between layers
\cite{dijkstra1982}.
